\subsection{Equivalence}

\begin{notation}
    \[P \lequiv Q \coloneqq (P \implies Q) \wedge (Q \implies P)\]
\end{notation}

% P=>Q, Q=>P > P<=>Q
\begin{align*}
    \premise{P \implies Q}\\
    \premise{Q \implies P}\\
    \byrule{(P \implies Q) \wedge (Q \implies P)}{\wedgeI}\\
    \bydefinition{P \lequiv Q}
\end{align*}
\newcommand{\lequivI}{{\lequiv}I}
\[\inferrule{P \implies Q \\ Q \implies P}{P \lequiv Q} \quad (\lequivI)\]

% P<=>Q > P=>Q, Q=>P
\begin{align*}
    \premise{P \lequiv Q}\\
    \bydefinition{(P \implies Q) \wedge (Q \implies P)}\\
    \byrule{P \implies Q}{\wedgeE}
\end{align*}
\begin{align*}
    \premise{P \lequiv Q}\\
    \bydefinition{(P \implies Q) \wedge (Q \implies P)}\\
    \byrule{Q \implies P}{\wedgeE}
\end{align*}
\newcommand{\lequivE}{{\lequiv}E}
\[\inferrule{P \lequiv Q}{P \implies Q} \quad \inferrule{P \lequiv Q}{Q \implies P} \quad (\lequivE)\]

% P/Q, P<=>Q > Q/P
\begin{align*}
    \premise{P}\\
    \premise{P \lequiv Q}\\
    \byrule{P \implies Q}{\lequivE}\\
    \byrule{Q}{\impliesE}\\
\end{align*}
\begin{align*}
    \premise{Q}\\
    \premise{P \lequiv Q}\\
    \byrule{Q \implies P}{\lequivE}\\
    \byrule{P}{\impliesE}\\
\end{align*}
\newcommand{\lequivE}{{\lequiv}E}
\[\inferrule{P \\ P \lequiv Q}{Q} \quad \inferrule{Q \\ P \lequiv Q}{P} \quad (\lequivE)\]

% > P<=>P
\begin{align*}
    \byrule{P \implies P}{\impliesR}\\
    \byrule{P \lequiv P}{\lequivI}
\end{align*}
\newcommand{\lequivR}{{\lequiv}R}
\[\vdash P \lequiv P \quad (\lequivR)\]

% P<=>Q > Q<=>P
\begin{align*}
    \premise{P \lequiv Q}\\
    \byrule{P \implies Q}{\lequivE}\\
    \byrule{Q \implies P}{\lequivE}\\
    \byrule{Q \lequiv P}{\lequivI}
\end{align*}
\newcommand{\lequivC}{{\lequiv}C}
\[\inferrule{P \lequiv Q}{Q \lequiv P} \quad (\lequivC)\]

% P<=>Q, Q<=>R > P<=>R
\begin{align*}
    \premise{P \lequiv Q}\\
    \premise{Q \lequiv R}\\
    \assumption[1]{P}\\
    \byrule[1]{Q}{\lequivE}\\
    \byrule[1]{R}{\lequivE}\\
    \byrule{P \implies R}{\impliesI}\\
    \byrule{R \lequiv Q}{\lequivC}\\
    \byrule{Q \lequiv P}{\lequivC}\\
    \assumption[1]{R}\\
    \byrule[1]{Q}{\lequivE}\\
    \byrule[1]{P}{\lequivE}\\
    \byrule{R \implies P}{\impliesI}\\
    \byrule{P \lequiv R}{\lequivI}
\end{align*}
\newcommand{\lequivA}{{\lequiv}A}
\[\inferrule{P \lequiv Q \\ Q \lequiv R}{P \lequiv R} \quad (\lequivA)\]
